\scanpage{26}
SYSTÈMES de RACINES (12) \hfill Tribu n° 59

§ 9. \uuline{Classification des \unsure{6tèmes} de racines.}

\uline{n° 1} - Dire qu'on est sur $R$, (l.2 du n° 1). A la l.5 bas de la p. 177 remplacer ``Il s'agit'' par ``Nous allons''. En bas de la p. 177, indexer par un ensemble $I$ à $\ell$ éléments, ce qui simplifie. Remonter la prop. 1 à la p. 178, (les murs étant indexés par $I$ et $I'$). Dessiner qq. chose dans le trou de la p. 178.

Il est antibourbachique que les lemmes aient des corollaires : les corollaires des pp. 179 sqq. seront qualifiés de lemmes. Dans le lemme 1, considérer l'ensemble $J$ des $j$ liés à $i$. Mieux expliciter le procédé d'identification de la p. 180 : on part d'une arête $(i,j)$ d'ordre 3 ; on enlève $i$, $j$ de l'ensemble des sommets, et on y met un nouveau sommet $p$ ; les arêtes sont, d'une part les anciennes ne contenant ni $i$ ni $j$, d'autre part, si $(i,k)$ ou $(j,k)$ est une vieille arête, $(p,k)$ sera une nouvelle arête ; les ordres sont conservés ; algébriquement, ça revient à remplacer les variables $\xi_i$ et $\xi_j$ par une nouvelle variable $\xi_p$.

Enoncé du haut de la p. 181, (devenu lemme) : ou bien un seul point de ramif. et toutes arêtes d'ordre 3, ou bien chaîne et au plus une arête d'ordre $\geq4$.

Les lemmes 3 et 4 ont l'air, soit canulés, soit d'énoncés tirés par les cheveux ; revenir à l'état 3, (n° 383), où c'était plus clair.

La prop. 1 remonte au début du n°.

\uline{n° 2} - Commencer par : nous allons voir que ces trois zèbres correspondent à des groupes ; y peuvent pas provenir de systèmes de racines à cause des 5 ou des $p\geq7$, (référence) ; par contre, nous verrons plus loin que tous les autres proviennent de systèmes de racines. A la fin du n°, dire que les formes obtenues

\mspage{25}
